


Constitutive modelling at finite strains is governed by two complementary and independent requirements. The first is \emph{thermodynamic consistency}, whereby constitutive equations must satisfy the laws of thermodynamics \cite{Gurtin2010,BonetGil2026}. The second is \emph{mathematical admissibility}, whereby the resulting boundary value problem remains well posed \cite{Ball_1976,Silhavy1997}. Neither requirement implies the other. Thermodynamically consistent constitutive models may still exhibit loss of ellipticity, unstable wave propagation or numerical breakdown, whereas mathematically stable formulations may violate fundamental thermodynamic principles. A robust constitutive theory must therefore enforce both requirements simultaneously \cite{Holzapfel2000}.

Within finite thermoelasticity and electro-mechanics, constitutive equations are typically derived to satisfy the First and Second Laws of Thermodynamics through an appropriate choice of Helmholtz free energy and enforcement of the Coleman-Noll procedure \cite{ColemanNoll1963}. This framework underpins most modern formulations of electro-elasticity, thermoelasticity and coupled multiphysics extensions at finite strains \cite{Miehe1994,HolzapfelSimo1996}. By contrast, the Third Law of Thermodynamics has received comparatively little attention in constitutive modelling \cite{Govindjee2024,Bonet2021first}. Recent work by Bonet and Gil \cite{BonetGil2026} has shown that a broad class of thermoelastic constitutive models, particularly those employing constant or purely temperature-dependent specific heat, are incompatible with the Nernst postulate, namely, \emph{the entropy of a system must approach zero as the temperature approaches absolute zero}. As a consequence, the Third Law introduces non-trivial restrictions on the admissible functional form. Specifically, strict adherence to the Third Law requires the use of deformation- and temperature-dependent specific heat and results in a fundamentally nonlinear stress–temperature coupling \cite{Bonet2021first,Simavilla2018}.

Mathematical admissibility introduces a second and independent set of constraints. In nonlinear elasticity, polyconvexity \cite{Ball_1976} —a condition stronger than rank-one convexity— ensures the existence of minimisers for the associated variational problems and plays a central role in guaranteeing numerical stability in nonlinear finite element implementations \cite{Bonet_book_statics,Bonet_polyconvexity_2015}. This property is critical not only for quasi-static simulations but also for dynamic scenarios where the preservation of hyperbolicity and stability of wave propagation is essential \cite{Bonet2021first}. These properties are particularly important in coupled multiphysics settings, where loss of ellipticity may be triggered by strong electromechanical or thermomechanical coupling. The extension of polyconvexity to coupled electro-magneto-mechanical formulations has been developed by Gil and Ortigosa [XXX], providing a consistent variational framework for multiphysics finite deformation problems. However, the incorporation of thermal restrictions associated with the Third Law remains unresolved.

Electro-Active Polymers (EAPs) constitute a particularly demanding class of materials in which these issues become critical [XXX]. They exhibit very large strains, low stiffness, and strong electromechanical coupling, which makes them attractive for applications involving large actuation and sensing capabilities [XXX]. In soft robotics, they have been widely explored as artificial muscles capable of replicating biological motion, enabling devices such as crawling, swimming and morphing structures. Their low weight, high compliance and multifunctionality have further motivated their use in biomedical devices, adaptive systems and soft actuators operating under large deformation regimes [XXX].

From a modelling perspective, EAPs represent a stringent benchmark for constitutive theories due to the coexistence of large strains, strong electromechanical coupling, thermal effects and time-dependent dissipation. Viscous effects and dielectric losses generate internal heat, leading to significant coupling between mechanical, electrical and thermal fields \cite{Reese_Govindjee_1998}. These effects become particularly relevant in dynamic actuation scenarios, where thermal accumulation can significantly alter both stiffness and dielectric response. The prediction of such behaviour therefore requires constitutive models that are not only thermodynamically consistent but also mathematically robust under extreme multiphysics coupling \cite{Bonet_polyconvexity_2015}. In this context, several contributions by the present authors and co-workers have developed polyconvex formulations for electro-mechanical, thermo-mechanical and electro-thermo-mechanical systems, establishing a variationally consistent framework for multiphysics finite elasticity \cite{Bonet_polyconvexity_2015,Bonet_cross_product_2016,Bonet2021first}. Nevertheless, existing formulations do not fully address the constraints imposed by the Third Law of Thermodynamics in combination with polyconvexity and finite-strain viscoelasticity.

The objective of the present work is to develop a constitutive framework for finite electro-thermo-viscoelasticity that simultaneously satisfies the three laws of thermodynamics and preserves stability. The formulation is based on a Helmholtz free-energy density depending on deformation, electric field, temperature and internal variables associated with the evolution of viscoelastic effects. The main contributions of the paper are threefold. First, a polyconvex electro-thermo-hyperelastic constitutive framework is developed that satisfies the First, Second and Third Laws of Thermodynamics while ensuring positive heat capacity and material stability under finite deformation. Second, a finite-strain viscoelastic extension is formulated through a consistent equilibrium/non-equilibrium decomposition that preserves thermodynamic structure and is directly extendable to anisotropic electro-active polymers via suitable structural tensors, exploiting a recent contribution by some of the authors \cite{Bonet_cross_product_2016}. Third, the proposed framework is calibrated against available thermo-mechanical, electro-mechanical and coupled electro-thermo-mechanical experiments and assessed through large-scale finite element simulations that quantify the role of viscous dissipation, heat generation and multiphysics coupling in the overall response.

The remainder of the paper is organised as follows. Section 2 introduces the governing equations of electro-thermo-mechanics and the associated thermodynamic restrictions. Section 3 develops the proposed constitutive framework. Section 4 presents material parameter identification. Section 5 demonstrates the performance of the model in representative finite element simulations. Finally, Section 6 summarises the principal conclusions and outlines possible extensions.


